% ------------------------
% example.tex
% ------------------------
% This file demonstrates various features of the note-thesis.cls class

\documentclass{NoteThesis/note-thesis}

\usepackage{xeCJK}
\setCJKmainfont{Noto Serif CJK SC}
\setCJKsansfont{Noto Sans CJK SC}
\setCJKmonofont{Noto Sans Mono CJK SC}
% ------------------------
% Document Information
% ------------------------

% Course name shown in the header (right side)
\course{Aesthetic Education}

% Title page information
\coursetitle{U30G11001 Aesthetic Education / 大学美育}  % Course title
\semester{2025B}                    % Academic semester
\teacher{Prof. Shuhan Yu}               % Instructor name

\begin{document}

% ------------------------
% Title Page
% ------------------------
% \maketitle generates a title page with course title, semester, and instructor
% A horizontal line is automatically drawn below
\maketitle

% ------------------------
% Table of Contents
% ------------------------
% \tableofcontents generates a two-column TOC with depth up to subsection (level 2)
\tableofcontents

\newpage


\section{审美需要}

审美需要是一种精神、情感的享受需要。在马斯洛需求理论中，审美需要在第六层。享受的需要往往由物质享受发展向求美的精神享受。

\subsubsection*{美育的功能}

\begin{enumerate}
    \item 培养审美能力
    \item 陶冶性情
    \item 完善人格
    \item 树立正确审美观
\end{enumerate}

\subsection{审美主体}

人有了审美需要就成为了真正的审美主体，是审美活动的发起者，承担者，在审美活动中起着\textbf{积极的主导作用}。

\subsubsection{\texorpdfstring{\mbox{\scalebox{0.75}{\fontspec{Noto Sans CJK SC}★}}}{} 审美主体的特点}

不同于实践主体、认识主体、伦理主体，具有自己的特征：

\begin{enumerate}
    \item 是\textbf{感性观照}的主体
    \item 是\textbf{情感活动}的主体
    \item 是\textbf{自由}的主体
\end{enumerate}

\subsection{审美欲望与审美兴趣}

从原发性的追求快乐的生命需要发展而来的层次，即审美欲望。表现为对形式、秩序的本能的追求。

审美欲望升华成为审美兴趣，即审美主体在审美活动中表现出的比较稳定的审美偏爱和倾向。

\subsection{审美理想}

是真善美的统一，从宏观方面确定链选择审美对象的范围，并体现了创造美的目标，激励人们努力创作美。

是人生境界的追求，也是审美追求的最高层次。

\section{审美能力}

审美心理因素包括：感觉、知觉、想象、情感、领悟。由此，审美能力包括：审美感知力，审美想象力，审美领悟力。

\subsection{审美感知力}

对质料的感受（生理性快感）构成了美感的一部分，但快感主要还是通过与其他更为高级的心理因素的联系参与美感的。两者的区别是：

\begin{enumerate}
    \item 美感是超实用、无功利的
    \item 美感的对象是审美意象，而快感没有
\end{enumerate}

\begin{mynote_env}[frametitle={知觉}]
    知觉是与感觉融合在感知中发生作用的，但知觉的对象是形式而非质料。
\end{mynote_env}

\subsection{审美想象力}

想象是人类心理感受对时间和空间的超越。

想象的简单形式是联想，包括接近联想、类比联想。更进一步地，想象具有更多主体自由性————生产新表象。包括：再造性想象、创造性相象。

其体现了：\textbf{自由性、创造性与情感性}，构成了独特的审美心理现象。

\begin{mynote_env}[frametitle={审美情感}]
    情感是审美和艺术的象征代表、甚至是特质所在。审美情感的对象是人的自由本质，也是人区别于动物的体验感受。
\end{mynote_env}

\begin{mynote_env}[frametitle={情景交融}]
    是中国古典美学概念，表现出审美的物我统一、主客同一的本质。
\end{mynote_env}

\subsection{审美领悟力}

审美领悟内在地贯穿审美感受全过程，把握了审美意蕴本质，与其他审美因素交融一体。

具有以下特点：

\begin{enumerate}
    \item \textbf{情感性}，普遍为参与性的情绪状态。
    \item \textbf{非概念化与感知性}，不是通过概念和逻辑，而是通过形象/意象表达某种本质性的东西。
    \item \textbf{趋向无穷性}，审美领域的对象不是有限性的实在物，而是在对象化中展示的人的自由创造性。
    \item \textbf{直觉性}，把握住审美意蕴的一刹那，就是审美感受的高峰，它使得审美领悟成为心理飞跃的关键节点，基本特征是直觉性。
\end{enumerate}

\subsection{\texorpdfstring{\mbox{\scalebox{0.75}{\fontspec{Noto Sans CJK SC}★}}}{} 审美经验模式}

需要后天培养，有五种经验模式：

\begin{enumerate}
    \item \textbf{审美体验}，亲身体验，源初性、亲历性、整体性，在审美经验中起基础性模式和奠基作用
    \item \textbf{审美注意}，客体吸引主体的注意力
    \item \textbf{审美移情}，由我及物，由物及我
    \item \textbf{完形心理}，欣赏整体大于部分
    \item \textbf{审美心理距离}，与现实拉开距离（当局与旁观），超越使用态度进行欣赏。
    \item \textbf{审美领悟}，如审美领悟力一节所述。
\end{enumerate}

\section{审美标准}

一般是在主体的审美活动中形成的，用于评价对象的一种内在尺度。审美标准主要有主体的内在尺度（个人的）、人类群体的标准（群体的一致性）和社会性标准（文化与道德）。

审美意识的客观标准的客观普遍性和历史具体性统一于特定的审美活动中，集中于时代性、民族性、阶级性三个反面。

\subsubsection*{审美标准的历史演变}

人的审美观会随着时代的变化而变化，由此形成不同的审美趣味，从而产生不同的审美经验。

纵向地看，美是随着时代发展而发展的。横向地看，不同文化间也有着不同的美，由民族、阶级和文化所塑造。

\subsection{\texorpdfstring{\mbox{\scalebox{0.75}{\fontspec{Noto Sans CJK SC}★}}}{} 形式美的基本法则}

\begin{enumerate}
    \item \textbf{对称与均衡}，对称两侧同质同量，具有稳定性；均衡是非对称的，两侧相互呼应，稳定中有动感，一致中有活泼。
    \item \textbf{调和与对比}，调和趋同、对比趋异；巧妙的调和/对比会引起对比、调和感。
    \item \textbf{比例与尺度}，主客体比例匀称，给人以舒适感。
    \item \textbf{节奏与韵律}，节奏是形式因素在运动中有规律的变化，韵律则是有变化的、岳动的规律感。
    \item \textbf{多样与统一}，事物个性的差别造就了多样，而统一是事物的共性与联系。
\end{enumerate}

\section{自然美}

总体上，存在外在自然物之美与内在天性之美两种内涵。相较于社会美、艺术美，自然美依赖于自然事物及其属性，内容较为宽泛朦胧而形式却清晰生动（偏向与形式美）。

\subsection{\texorpdfstring{\mbox{\scalebox{0.75}{\fontspec{Noto Sans CJK SC}★}}}{} 自然美的代表性模式}

\begin{enumerate}
    \item \textbf{如画模式}，古典烂漫
    \item \textbf{比德模式}，儒家人格化的道德美
    \item \textbf{宇宙模式}，毕达格拉斯式
    \item \textbf{天成之美}，自然无为的内在天成之美
\end{enumerate}

\section{社会美}

以人类最基础的生存活动————劳动为核心而形成的审美形态，是人类最早最基本的审美领域之一，包括：

\begin{enumerate}
    \item \textbf{社会实践活动的美}，发现于生活中
    \item \textbf{社会实践成果的美}，包括人文环境、人类文化遗存、文化传统等
    \item \textbf{社会实践主体的美}，人的美
\end{enumerate}

\subsection{社会美的特征}

\begin{enumerate}
    \item 是人类社会实践产物最直接的现实存在形式
    \item 直接体现了人的自由自觉的特性
    \item 体现出明显的价值取向
    \item 形式感从属于其美的内容
\end{enumerate}

\subsection{\texorpdfstring{\mbox{\scalebox{0.75}{\fontspec{Noto Sans CJK SC}★}}}{} 人物美}

\begin{enumerate}
    \item \textbf{人体美}，相貌、仪表等方面的美，自然与社会因素的统一，最终显示出了人的感性生命。
    \item \textbf{人格美}，人形体、行为中体现出的独特人格魅力与崇高精神境界。人对自身精神力量的崇敬和向往。
\end{enumerate}

\begin{mynote_env}[frametitle={节日欢庆的美}]
    超越日常生活的刻版与功利，回归原始本真的存在形态。酒神而非日神式的狂欢。
\end{mynote_env}

\subsection{技术美与功能美}

技术美为，在技术领域有审美价值、有明确使用目的和功能。

功能美为，从人的作为使用对象的角度分析技术美。

\section{\texorpdfstring{\mbox{\scalebox{0.75}{\fontspec{Noto Sans CJK SC}★}}}{} 艺术美}

艺术作品是艺术家创造性劳动的产物，是一种特殊的精神产品。艺术作品的美，即艺术作品的审美价值。

\subsection{中外对艺术的共同认识}

\begin{enumerate}
    \item \textbf{人工性}，艺术必然存在一种物质载体，即艺术品。
    \item \textbf{技艺性}，艺术性与技艺与才能密不可分。
    \item \textbf{精神性}，不止为了使用，也是为了满足人的精神需求。
\end{enumerate}

\subsection{艺术本体的主要观点}

\begin{enumerate}
    \item \textbf{模仿说}，是现实世界的模仿描摹与再现。
    \item \textbf{表现说}，艺术是情感，内心世界和主观精神的表现。
    \item \textbf{形式说}，艺术的本体在于形式或纯粹的形式。
    \item \textbf{惯例说}，艺术是由一定时代人们的习俗规定的————换言之，艺术是由艺术界公众决定的。
\end{enumerate}

从中国美学总结，艺术的本体是审美意象，“情”、“景”统一，呈现一个完整的有意蕴的感性世界。

\subsection{艺术作品的层次结构}

\begin{enumerate}
    \item \textbf{材料层}，影响意象的生成，给观赏者一种质料感。
    \item \textbf{形式层}，是材料的形式化，体现了作品（亦即整个意象世界）的意味和形式本身的意味。
    \item \textbf{意蕴层}，「意蕴」是美感的对象，带有多义性，是一个永无止境的历史显现的过程和永无止境的生成的过程。
\end{enumerate}

\subsection{意境}

所谓意境，就是超越具体的、有限的物象、事件、场景，进入无穷的时间与空间，获得一种哲理性的感悟。

意境即康德所谓「惆怅」，尼采所谓「形而上的慰藉」，不止存在于艺术美中，也存在与社会美和自然美领域中。

\section{\texorpdfstring{\mbox{\scalebox{0.75}{\fontspec{Noto Sans CJK SC}★}}}{} 审美本质}

\subsection{西方美学史上的审美本质}

\begin{enumerate}
    \item \textbf{柏拉图}，美的现象和事物背后的「美」是理念
    \item \textbf{亚里士多德}，美的本质在于感性事物的本身，在于事物的形式和比例
    \item \textbf{毕达哥拉斯}，美是和谐与比例
    \item \textbf{休谟}，美是主观的心理经验，「快乐」的感觉
    \item \textbf{康德}，美是一对象的合目的性的形式（审美心理结构论）
    \item \textbf{黑格尔}，美是感性与理性的统一，内容与形式的统一，是理念的感性显现
\end{enumerate}

在现当代，西方美学研究也认为，美是一种直觉（克罗齐），美由「内模仿」产生（谷鲁斯、浮龙 李），或是「移情」产生（立普斯），距离产生美（布洛）。

\subsection{马克思「美的规律」思想}

审美是按照「美的规律」构造的、体现人的本质的自由劳动。

「美的规律」、「自由劳动」、「人的本质」三个概念三位一体。

\subsection{中国美学史上的审美本质}

美是人内在修养的自然呈现

\subsubsection*{儒家美学：以人品为中心}

\begin{enumerate}
    \item \textbf{孔子，以和为美}，「美」与「善」、「文」与「质」的统一
    \item \textbf{孟子，充实为美}，充实而有光辉之谓之大
\end{enumerate}

\subsubsection*{道家美学：大象无形}

强调「道」对审美表现的本体意义。

\section{\texorpdfstring{\mbox{\scalebox{0.75}{\fontspec{Noto Sans CJK SC}★}}}{} 优美与崇高}

\subsection{优美}

优美，也称阴柔之美，具有小巧、平静、舒缓和单纯等表现形式，\textbf{和谐、完满}的典型特征，给人以和谐、完满的审美。

\subsubsection{优美的本质}

优美是人的本质的对象化的\textbf{实现和结构}，是合目的性与合规律性的统一。

优美是人本质对象化的活动结果的展示，在本质上是静态的，很少有矛盾冲突和审美张力的展示。具有柔和平稳、舒畅喜悦的美感心理特质。

\subsection{崇高}

崇高是一种突出了主体与客体，人与自然，感性与理性的矛盾的复杂心理体验，以\textbf{痛感、压抑}为基础，突出\textbf{不和谐到和谐、痛感到快感的复杂体验}。具有狂放，暴烈，无限与模糊的特征。

崇高的力量根植于个体向总体性人类的更高的升华，即主体精神的张扬和无限延伸（特征），侧重于社会美和艺术美的展示，在庄严宏伟的气势中表示主体斗争的痕迹和强烈的伦理道德力量。

\subsubsection{崇高的本质}

崇高展示人的本质对象化实现的\textbf{曲折过程}，展示主客体从巨大冲突到趋向统一的艰难过程，借此推动实践主体的精神力量的生成和提升。

\subsection{崇高与优美的转化}

审美范畴的转化本质是主客体对立统一的矛盾运动，主体的地位的变化形成了如下序列

\begin{center}
    崇高-悲剧-喜剧-优美
\end{center}

优美往往是审美范畴发展的基本趋向。

\begin{mynote_env}[frametitle={现代与后现代的审美范畴}]
    \textbf{现代}艺术的突出变化是崇高的解构与丑的出现，即「不再崇高的英雄和不再美丽的艺术」。\\
    \textbf{后现代}的审美贫困，无疑与主体的变异、沉沦，及崇高的历史意义与目标的模糊存在关联。真、善之间的冲突与错位，目的性和善间关联的消解，导致主体向唯我主体滑落并产生文化上的虚无主义。
\end{mynote_env}

\section{悲剧与喜剧}

在不同语境下，悲剧与喜剧不止作为喜剧体裁的形式存在，也作为审美范畴或审美类型的分类而存在，并描述一种审美体验和评价。

\subsection{悲剧}

因此，悲剧也可以被解读为：

\begin{enumerate}
    \item \textbf{日常生活}，不幸或悲惨的事情，作为创作素材和现实基础存在
    \item \textbf{戏剧体裁}，相对于喜剧和正剧的体裁，作为典型代表
    \item \textbf{审美范畴}，与喜剧、优美、崇高并列，代表哲学概括
\end{enumerate}

现实中的苦难往往先是「需要严肃对待的伦理事件」，因此只有拉开心理距离，悲惨事件才能作为审美形态的悲剧呈现出来。

\subsubsection{\texorpdfstring{\mbox{\scalebox{0.75}{\fontspec{Noto Sans CJK SC}★}}}{} 悲剧的美学范畴}

悲剧美，指具有值得同情和认同的个体，在不可避免的社会冲突中，遭遇不应有又无法逃避的不幸结局，个性遭到毁灭，激发审美者的悲伤、怜悯和恐惧等复杂情感的审美形态

\begin{mynote_env}[frametitle={悲剧的本质}]
    悲剧的本质在于对个性或人性毁灭事件自觉或不自觉的哲学性观照
\end{mynote_env}

\subsubsection{悲剧的特征}

\begin{enumerate}
    \item 悲剧至少有一个值得同情的人物作为主人公
    \item 主人公在不可避免的社会冲突中，遭遇不幸结局，个性遭到毁灭
    \item 不幸的结局是无法逃避，但又不正当、令人无法接受的。
    \item 会引发审美者的恐惧、悲伤或怜悯的情感和审美感受，甚至发生某些精神提升或转变。
\end{enumerate}

相较于崇高，悲剧更倾向于不可避免的失败而非成功前景，审美对象的范围更小，只包括人的行为及其事件。

\subsection{喜剧}

喜剧是存在各种弱点、缺陷甚至虚假丑恶的事物，在特定的冲突或情景下暴露出自身的不协调与自相矛盾性，从而令人发笑的审美形态。

\begin{mynote_env}[frametitle={喜剧的类型}]
    喜剧由对主人公的价值态度可以分为「肯定性喜剧」和「否定性喜剧」（或讽刺性、批判性）。
\end{mynote_env}

\begin{mynote_env}[frametitle={喜剧的本质}]
    喜剧的本质是对作为审美对象的主人公遮蔽自我个性或人性本质的事件的揭露。
\end{mynote_env}

\subsubsection{喜剧的特征}

\begin{enumerate}
    \item 喜剧人物的无价值或反价值
    \item 喜剧事件的矛盾与可笑性
    \item 矛盾的突发性
    \item 喜剧效果的愉悦性（喜剧的愉悦性来得更为直接和大众性）
\end{enumerate}

\subsubsection*{喜作为审美价值}

喜的价值载体往往是被历史实践否定了的，与正常尺度相偏离了的事物。

喜的对象实质是被审美主体把握的。

\section{丑}

丑是一种特殊的否定性的审美范畴，广义上的丑是现实对人本质的否定，是合目的性与和规律性的背离，狭义上的丑是与优美相对立，是对其纯粹、完满和和谐性的否定。

\subsection{\texorpdfstring{\mbox{\scalebox{0.75}{\fontspec{Noto Sans CJK SC}★}}}{} 美丑共存中丑的特征}

在生活中有美丑共存的复杂局面，在审美活动中必将包括对美的欣赏和创造、与对丑的揭露与摒弃，因此，美与丑都不是作为物化形态的结果存在于现实世界中的，而是在审美实践中不断生产和发展的。

在特定情况下，丑可以让人联想到、或转化为美。

\subsection{\texorpdfstring{\mbox{\scalebox{0.75}{\fontspec{Noto Sans CJK SC}★}}}{} 丑的审美价值}

\begin{enumerate}
    \item \textbf{对比和陪衬美的价值}，增大心理的反差，增强美的光辉
    \item \textbf{丰富和激活美的价值}，狭义的丑进入审美领域，以其混乱不和谐特征成为形式美中的特殊形态，扩宽了美的内涵。
\end{enumerate}

\section{荒诞}

荒诞作为审美形态和范畴是在现代确立的，它建立在现代社会的高度异化和人对异化社会的绝望和孤独感上。

荒诞是对无意义、无本质、无价值之本身的对象化，因此也是对象化的虚无和虚无的对象化。

\subsection{荒诞的审美价值}

\begin{enumerate}
    \item \textbf{形式的碎片化和丑怪}，荒诞不再将实体看作可以理性透彻了解的存在
    \item \textbf{内容的平面化与无中心}，内容的中心和主体弥散，出现了无戏剧高潮和情节的倾向
    \item \textbf{主体的零度化和无力感}，主体不再如传统审美形态中敢于并有能力做出并表达价值判断————前景的对象都是等值或无价值的。
\end{enumerate}

荒诞的行为正是对荒诞的反抗，这反抗是为了人从荒诞中获得幸福。

\end{document}
